

\section{Conclusions}

The project demonstrates that hybrid energy systems can significantly reduce carbon emissions in the studied area by fully eliminating fossil 
fuel use for electricity generation. Moreover, it lays the groundwork for further decarbonization of the transportation and heating sector.

The plant model confirms the technical feasibility of the proposed system, providing a detailed, hour-by-hour account of energy and material 
flows both within the plant boundaries and between the plant and the surrounding community.

The plant layout was, at first, roughly designed primarily around the operational requirements of the gas turbine. As a result, the number of High-Temperature 
Steam Electrolyzers (HTSEs) and Small Modular Reactors (SMRs) was determined based on this component constraints. 
The plant was accordingly oversized, in order to ensure sufficient syngas production to let the gas turbine operate at full capacity.

The ambitious decarbonization goals led to a substantial portion of the nuclear-generated electricity being allocated to 
low-profit decarbonization processes — such as hydrogen and methane production — rather than being sold directly to the grid.

A key finding of the economic analysis is that the high initial capital costs of the full hybrid system are unlikely to be recovered if all 
components are built simultaneously. However, a phased implementation strategy can mitigate these costs and ultimately generate positive returns.

To manage the financial burden and benefit from advancements in relevant technologies, the installation of the 
HTSE-Methanator-Gas Turbine complex is postponed until the 20th year of operation. 
In the first phase, the electricity from the SMRs is sold directly to the grid, generating more immediate and profitable returns 
to help offset nuclear capital expenditures.

This phased strategy mirrors the real-world approach of deploying SMRs incrementally to finance subsequent units. 
Moreover, this strategic choice demonstrates resilience with respect to the main financial parameters possibly affecting its economic viability.

When combined with the project's strong decarbonization potential and the elimination of foreign energy dependence, these results make 
the proposal a compelling opportunity for the Seven Lakes region and a strategic use of the JRC's nuclear infrastructure.